\chapter{The Buffer Holds Facts}\label{record-036}

\noindent{\footnotesize\color{BookMuted}\path{docs/invariants/buffer-holds-facts.md}\par}\smallskip

Buffers, queues, and message channels should carry facts about what happened or what should happen next. They should not carry executable behavior.

\section{Meaning}\label{record-036-meaning}

Good buffered values look like:

\begin{itemize}
\tightlist
\item
  \texttt{\{\allowbreak{} \allowbreak{}t\allowbreak{}y\allowbreak{}p\allowbreak{}e\allowbreak{}:\allowbreak{} \allowbreak{}'\allowbreak{}q\allowbreak{}u\allowbreak{}i\allowbreak{}t\allowbreak{}-\allowbreak{}r\allowbreak{}e\allowbreak{}q\allowbreak{}u\allowbreak{}e\allowbreak{}s\allowbreak{}t\allowbreak{}e\allowbreak{}d\allowbreak{}'\allowbreak{} \allowbreak{}\}}
\item
  \texttt{\{\allowbreak{} \allowbreak{}t\allowbreak{}y\allowbreak{}p\allowbreak{}e\allowbreak{}:\allowbreak{} \allowbreak{}'\allowbreak{}n\allowbreak{}o\allowbreak{}t\allowbreak{}i\allowbreak{}f\allowbreak{}i\allowbreak{}c\allowbreak{}a\allowbreak{}t\allowbreak{}i\allowbreak{}o\allowbreak{}n\allowbreak{}-\allowbreak{}a\allowbreak{}d\allowbreak{}d\allowbreak{}e\allowbreak{}d\allowbreak{}'\allowbreak{},\allowbreak{} \allowbreak{}l\allowbreak{}e\allowbreak{}v\allowbreak{}e\allowbreak{}l\allowbreak{}:\allowbreak{} \allowbreak{}'\allowbreak{}i\allowbreak{}n\allowbreak{}f\allowbreak{}o\allowbreak{}'\allowbreak{},\allowbreak{} \allowbreak{}t\allowbreak{}e\allowbreak{}x\allowbreak{}t\allowbreak{}:\allowbreak{} \allowbreak{}'\allowbreak{}S\allowbreak{}a\allowbreak{}v\allowbreak{}e\allowbreak{}d\allowbreak{}'\allowbreak{} \allowbreak{}\}}
\item
  \texttt{\{\allowbreak{} \allowbreak{}t\allowbreak{}y\allowbreak{}p\allowbreak{}e\allowbreak{}:\allowbreak{} \allowbreak{}'\allowbreak{}f\allowbreak{}r\allowbreak{}a\allowbreak{}m\allowbreak{}e\allowbreak{}-\allowbreak{}a\allowbreak{}c\allowbreak{}t\allowbreak{}i\allowbreak{}o\allowbreak{}n\allowbreak{}'\allowbreak{},\allowbreak{} \allowbreak{}a\allowbreak{}c\allowbreak{}t\allowbreak{}i\allowbreak{}o\allowbreak{}n\allowbreak{}:\allowbreak{} \allowbreak{}'\allowbreak{}p\allowbreak{}u\allowbreak{}s\allowbreak{}h\allowbreak{}-\allowbreak{}n\allowbreak{}o\allowbreak{}t\allowbreak{}i\allowbreak{}f\allowbreak{}i\allowbreak{}c\allowbreak{}a\allowbreak{}t\allowbreak{}i\allowbreak{}o\allowbreak{}n\allowbreak{}'\allowbreak{},\allowbreak{} \allowbreak{}n\allowbreak{}o\allowbreak{}t\allowbreak{}i\allowbreak{}f\allowbreak{}i\allowbreak{}c\allowbreak{}a\allowbreak{}t\allowbreak{}i\allowbreak{}o\allowbreak{}n\allowbreak{} \allowbreak{}\}}
\end{itemize}

Bad buffered values look like:

\begin{itemize}
\tightlist
\item
  class instances with \texttt{a\allowbreak{}p\allowbreak{}p\allowbreak{}l\allowbreak{}y\allowbreak{}(\allowbreak{})} or \texttt{r\allowbreak{}u\allowbreak{}n\allowbreak{}(\allowbreak{})} methods
\item
  closures that capture ambient state and perform side effects later
\item
  ``smart'' objects whose meaning only exists when executed
\end{itemize}

The rule is simple:

\begin{quote}
the buffer should hold facts, not code
\end{quote}

Interpretation belongs at the boundary that owns side effects.

\section{Why}\label{record-036-why}

Fact-shaped buffered values are:

\begin{itemize}
\tightlist
\item
  serializable
\item
  inspectable
\item
  loggable
\item
  replayable
\item
  testable without executing host effects
\end{itemize}

Behavior-shaped buffered values make the system harder to audit because the meaning of the queue is hidden inside arbitrary executable code.

\section{Applies to}\label{record-036-applies-to}

This invariant applies to:

\begin{itemize}
\tightlist
\item
  command buffers
\item
  event channels
\item
  queued frame actions
\item
  cross-thread messages
\item
  persisted or replayable runtime logs
\end{itemize}

It is especially important anywhere we may later want:

\begin{itemize}
\tightlist
\item
  diagnostics
\item
  crash replay
\item
  deterministic tests
\item
  worker transport
\item
  post-mortem inspection
\end{itemize}

\section{Does not mean}\label{record-036-does-not-mean}

This invariant does \textbf{not} mean every internal helper must become data.

It does not apply to:

\begin{itemize}
\tightlist
\item
  direct hot-path loops where no queue/buffer boundary exists
\item
  local strategy tables used only as implementation detail
\item
  pure helper functions that never become buffered values
\end{itemize}

The rule is about what crosses a buffered or queued boundary, not about banning functions from the codebase.

\section{Review questions}\label{record-036-review-questions}

When adding a queue or buffered protocol, ask:

\begin{enumerate}
\def\labelenumi{\arabic{enumi}.}
\tightlist
\item
  Could I print every buffered item as plain data and understand the system?
\item
  Could I serialize and replay this later?
\item
  Is the side-effect owner explicit and centralized?
\item
  Would a new engineer have to execute the object to learn what it means?
\end{enumerate}

If the answer to the last question is yes, the design is probably violating this invariant.

\section{Protected by}\label{record-036-protected-by}

\begin{itemize}
\tightlist
\item
  \href{https://github.com/flyingrobots/bijou/blob/1853ad902f6bf219326409e7f743a5a1d6517dbc/docs/legends/RE-runtime-engine.md}{RE --- Runtime Engine}
\item
  \href{https://github.com/flyingrobots/bijou/blob/1853ad902f6bf219326409e7f743a5a1d6517dbc/docs/legends/DX-developer-experience.md}{DX --- Developer Experience}
\end{itemize}

